\subsection*{期值、\texttt{.data} 与 \texttt{.span}}

DMA 传输可以返回一个期值句柄：

\begin{lstlisting}[language=croqtile,caption={由一次拷贝得到的期值}]
f = dma.copy input => shared;
\end{lstlisting}

它提供两个成员：

\begin{itemize}[leftmargin=1.6em]
\item \co@f.data@：传输的目的缓冲区。
\item \co@f.span@：目的缓冲区的形状。
\end{itemize}

使用方式如下：

\begin{lstlisting}[language=croqtile,caption={读取 data 与 span}]
f = dma.copy input => shared;
foreach idx in [f.span]
  output.at(idx) = f.data.at(idx);
\end{lstlisting}

即使拷贝是同步的，给结果命名也有用途：后续代码既能通过 \co@.data@ 访问目的缓冲区，又能通过 \co@.span@ 确定遍历范围。这样不必再单独维护一份缓冲区名称和形状。

图~\ref{fig:f04-fig_future_data} 展示了拷贝返回的句柄暴露了哪些信息。

\begin{figure}[htbp]\centering
\includegraphics[width=\linewidth,height=0.56\textheight,keepaspectratio]{f04-fig_future_data_light.png}
\caption{复制结果通过 .data 提供目标张量，通过 .span 提供其形状；异步复制必须完成后才能读取。}
\label{fig:f04-fig_future_data}
\end{figure}
